\documentclass[12pt]{article}
\usepackage[utf8]{inputenc}
\usepackage[T1]{fontenc}
\usepackage[margin=1in]{geometry}
\usepackage{setspace}
\usepackage{xcolor}
\usepackage[colorlinks=true, urlcolor=blue!60!black, linkcolor=blue!60!black]{hyperref}
\setlength{\parindent}{0em}
\setlength{\parskip}{0.45em}
\renewcommand{\baselinestretch}{1.0}
\begin{document}
\begin{center}
{\Large \textit{INS v. Chadha}, 462 U.S. 919 (1983)\par}
\vspace{0.3em}
{\normalsize Case Brief --- One-Page Edition\par}
\end{center}

\textbf{Memory Jogger.} The legislative veto case: bicameralism and presentment, no shortcuts (\href{https://supreme.justia.com/cases/federal/us/462/919/}{462 U.S. 919}).

\textbf{Facts.} An Immigration Judge suspended the deportation of Jagdish Rai Chadha, who had overstayed a student visa. Under \S~244(c)(2) of the immigration statute, the House alone, without debate or a recorded vote, passed a resolution disapproving the suspension; it never went to the Senate or the President.

\textbf{Procedural History.} Chadha was ordered deported; the Ninth Circuit held the veto unconstitutional.

\textbf{Votes.} Affirmed, 7--2. Burger wrote for the Court; Powell concurred in the judgment; White and Rehnquist dissented.

\textbf{Holding.} The one-House veto is legislative action that violates the bicameralism and presentment requirements of Article I, \S\S~1 and 7; it is severable from the rest of the statute.

\textbf{Reasoning.} Action that alters the legal rights of persons outside the Legislative Branch is legislative, and the Constitution's express unicameral exceptions foreclose others; Congress must abide by its delegation until it legislates. Powell saw a judicial act; White defended the veto as an ``indispensable, political invention'' of the administrative state.

\textbf{Examples.} \emph{Same-Side:} a two-House veto of agency rules. \emph{Opposite-Side:} a report-and-wait law with disapproval by joint resolution. \emph{Fence-Sitter:} required joint approval before major rules take effect.

\textbf{Critique.} Footnote 16 treats agency rulemaking as executive but congressional review as legislative, though both alter outsiders' rights, and severance enlarged executive power.

\textbf{Key Quotation.} ``It emerges clearly that the prescription for legislative action in Art. I, \S\S~1, 7, represents the Framers' decision that the legislative power of the Federal Government be exercised in accord with a single, finely wrought and exhaustively considered, procedure.'' --- Chief Justice Burger, Opinion of the Court, 462 U.S. at 951.

\end{document}
